% AC: abstract written last night, numbers are from FINAL_TEST.md, pls dont change them without checking the metrics json
\begin{abstract}
Open-vocabulary detectors like \gdino can detect any class that is written in the text prompt, but on aerial images giving the detector a long list of fine-grained and overlapping class names make the accuracy collapse almost to zero.
In this work we ask a simple question, can we choose a small vocabulary for every image \emph{before} detection, without any training, so that a frozen detector works better?
We propose a training-free selector that score all 80 candidate classes with the remote sensing CLIP model RemoteCLIP and keeps every class whose probability is above a threshold calibrated once on held-out images, so the size of vocabulary adapt to each image (around 8 classes on average).
The selector only see the pixels and cost about 8\,ms per image.
We evaluate on \lae with a source-aware protocol that account for its five disjoint label spaces, all our choices were frozen before a single run on 2,862 held-out test images.
The selector raise AP50 of \gdino-T from 0.0053 to 0.0280 under standard COCO scoring and from 0.0157 to 0.0418 under source-aware scoring, both robust under three statistical checks, and it is also 1.21$\times$ faster end-to-end.
However the gain is mostly on imagery similar to the encoder's training data (DIOR/DOTA, and NWPU VHR-10 without retuning) and it disappear on FAIR1M and xView.
\TODO{one more line on negative results? maybe drop if over length}
\end{abstract}
